\documentclass[12pt,a4paper]{article}

% Пакеты для русского языка
\usepackage[utf8]{inputenc}
\usepackage[russian]{babel}
\usepackage[T2A]{fontenc}

% Пакеты для оформления
\usepackage{geometry}
\usepackage{graphicx}
\usepackage{setspace}
\usepackage{indentfirst}
\usepackage{hyperref}
\usepackage{fancyhdr}
\usepackage{lipsum}
\usepackage{xcolor}
\usepackage{tikz}
\usepackage{subcaption}
\usepackage{wrapfig}
\usepackage{float}

% Настройка полей страницы
\geometry{
    left=30mm,
    right=15mm,
    top=20mm,
    bottom=20mm
}

% Настройка межстрочного интервала
\onehalfspacing

% Настройка отступа первой строки
\setlength{\parindent}{1.25cm}

% Настройка колонтитулов
\pagestyle{fancy}
\fancyhf{}
\fancyhead[R]{\thepage}
\fancyhead[L]{\leftmark}
\renewcommand{\headrulewidth}{0.4pt}

% Настройка гиперссылок
\hypersetup{
    colorlinks=true,
    linkcolor=blue,
    urlcolor=blue,
    citecolor=black
}

\begin{document}

% Титульная страница
\begin{titlepage}
    \centering
    
    \vspace*{1cm}
    
    {\large МИНИСТЕРСТВО НАУКИ И ВЫСШЕГО ОБРАЗОВАНИЯ\\
    РОССИЙСКОЙ ФЕДЕРАЦИИ}
    
    \vspace{0.5cm}
    
    {\large Федеральное государственное автономное образовательное учреждение\\
    высшего образования}
    
    \vspace{0.3cm}
    
    {\large\bfseries «Российский Университет Дружбы народов имени Патриса Лумумбы»}
    
    \vspace{0.5cm}
    
    {\large Факультет физико-математических и естественных наук}
    
    {\large Кафедра теории вероятностей и кибербезопасности}
    
    \vspace{3cm}
    
    {\Large\bfseries ОТЧЕТ}
    
    \vspace{0.3cm}
    
    {\large по лабораторной работе №4}
    
    \vspace{0.5cm}
    
    {\large на тему:}
    
    \vspace{0.3cm}
    
    {\large\bfseries «Работа с изображениями, плавающими объектами\\и перекрестными ссылками в LaTeX»}
    
    \vfill
    
    \begin{flushright}
        Выполнил:\\
        студент группы НПМмд02-24\\
        Викторов Егор Игоревич\\
        \vspace{0.5cm}
    \end{flushright}
    
    \vfill
    
    {\large Москва, 2025}
    
\end{titlepage}

% Содержание
\tableofcontents
\newpage

% Основная часть
\section{Введение}

Данная лабораторная работа посвящена изучению возможностей системы LaTeX по работе с графическими элементами, плавающими объектами и системой перекрестных ссылок. 

\subsection{Цель работы}

Целью данной работы является освоение практических навыков работы с изображениями в LaTeX, изучение параметров позиционирования графики и механизма перекрестных ссылок.

\subsection{Задачи работы}

В рамках лабораторной работы необходимо выполнить следующие задачи:

\begin{enumerate}
    \item Создать собственные изображения с использованием пакета TikZ
    \item Исследовать параметры width, height, angle и scale для изображений
    \item Изучить разницу между textwidth и linewidth
    \item Протестировать поведение изображений в режиме twocolumn
    \item Исследовать спецификаторы позиционирования плавающих объектов
    \item Изучить работу команд label и ref
    \item Определить количество компиляций для корректной работы ссылок
    \item Исследовать правильное размещение label относительно caption
\end{enumerate}

\section{Теоретическая часть}

\subsection{Пакет graphicx}

Пакет \texttt{graphicx} является стандартным инструментом для включения изображений в документы LaTeX. Основная команда для вставки изображений:

\begin{verbatim}
\includegraphics[параметры]{имя_файла}
\end{verbatim}

\subsection{Параметры изображений}

Основные параметры команды \texttt{includegraphics}:

\begin{itemize}
    \item \textbf{width} — ширина изображения
    \item \textbf{height} — высота изображения
    \item \textbf{scale} — масштаб изображения
    \item \textbf{angle} — угол поворота изображения
    \item \textbf{keepaspectratio} — сохранение пропорций
\end{itemize}

\subsection{Плавающие объекты}
Плавающие объекты (floats) в LaTeX позволяют системе автоматически размещать изображения и таблицы в оптимальных местах документа. Спецификаторы позиции:

\begin{table}[h]
\centering
\caption{Спецификаторы позиционирования}
\label{tab:specifiers}
\begin{tabular}{|c|l|}
\hline
\textbf{Спецификатор} & \textbf{Описание} \\
\hline
h & here — примерно в месте вставки \\
\hline
t & top — вверху страницы \\
\hline
b & bottom — внизу страницы \\
\hline
p & page — на отдельной странице \\
\hline
! & игнорировать некоторые ограничения \\
\hline
H & строго здесь (требует пакет float) \\
\hline
\end{tabular}
\end{table}

\section{Практическая часть}

\subsection{Создание собственных изображений с TikZ}

Для демонстрации возможностей LaTeX создадим несколько графических элементов с использованием пакета TikZ.

\subsubsection{Простая геометрическая фигура}

На рисунке \ref{fig:tikz1} представлена простая геометрическая композиция, созданная средствами TikZ.

\begin{figure}[h]
\centering
\begin{tikzpicture}
    % Прямоугольник
    \draw[fill=blue!30, draw=blue!50!black, line width=2pt] (0,0) rectangle (4,3);
    % Круг
    \draw[fill=red!30, draw=red!50!black, line width=2pt] (2,1.5) circle (1cm);
    % Треугольник
    \draw[fill=green!30, draw=green!50!black, line width=2pt] 
        (1,2.5) -- (3,2.5) -- (2,4) -- cycle;
    % Текст
    \node at (2,0.5) {\Large\textbf{TikZ}};
\end{tikzpicture}
\caption{Геометрическая композиция, созданная в TikZ}
\label{fig:tikz1}
\end{figure}

Как видно из рисунка \ref{fig:tikz1}, TikZ позволяет создавать векторную графику непосредственно в коде LaTeX.

\subsubsection{Диаграмма с градиентом}

Рисунок \ref{fig:tikz2} демонстрирует более сложную графику с использованием градиентов и узлов.

\begin{figure}[htbp]
\centering
\begin{tikzpicture}[scale=1.2]
    % Оси координат
    \draw[->, thick] (-0.5,0) -- (5,0) node[right] {$x$};
    \draw[->, thick] (0,-0.5) -- (0,4) node[above] {$y$};
    
    % График функции
    \draw[domain=0:4.5, smooth, variable=\x, blue, line width=1.5pt] 
        plot ({\x}, {0.5*\x*\x - 0.3*\x + 0.5});
    
    % Точки
    \fill[red] (1,0.7) circle (2pt) node[above right] {$A$};
    \fill[red] (2,1.9) circle (2pt) node[above right] {$B$};
    \fill[red] (3,4.1) circle (2pt) node[above right] {$C$};
    
    % Сетка
    \draw[gray!30, very thin] (0,0) grid (4.5,4);
\end{tikzpicture}
\caption{График параболической функции}
\label{fig:tikz2}
\end{figure}

\subsection{Исследование параметров width, height, angle и scale}

\subsubsection{Параметр width}

Создадим изображение и будем изменять его ширину относительно \texttt{textwidth}.

\begin{figure}[h]
\centering
\begin{tikzpicture}
    \draw[fill=orange!40] (0,0) rectangle (3,2);
    \node at (1.5,1) {\textbf{Оригинал}};
\end{tikzpicture}
\caption{Исходное изображение для экспериментов}
\label{fig:original}
\end{figure}

Теперь применим различные значения ширины:

\begin{figure}[H]
\centering
\begin{subfigure}{0.3\textwidth}
    \centering
    \begin{tikzpicture}[scale=0.5]
        \draw[fill=orange!40] (0,0) rectangle (3,2);
        \node at (1.5,1) {\tiny\textbf{0.3tw}};
    \end{tikzpicture}
    \caption{width=0.3\textbackslash textwidth}
\end{subfigure}
\hfill
\begin{subfigure}{0.45\textwidth}
    \centering
    \begin{tikzpicture}[scale=0.75]
        \draw[fill=orange!40] (0,0) rectangle (3,2);
        \node at (1.5,1) {\small\textbf{0.45tw}};
    \end{tikzpicture}
    \caption{width=0.45\textbackslash textwidth}
\end{subfigure}
\caption{Сравнение различных значений width}
\label{fig:width_comparison}
\end{figure}

\subsubsection{Параметр angle}

Рассмотрим поворот изображения на различные углы (рисунок \ref{fig:angles}).

\begin{figure}[htbp]
\centering
\begin{tikzpicture}
    % 0 градусов
    \begin{scope}[xshift=0cm]
        \draw[fill=cyan!30] (0,0) rectangle (1.5,1);
        \node at (0.75,0.5) {$0°$};
    \end{scope}
    
    % 45 градусов
    \begin{scope}[xshift=3cm, rotate=45]
        \draw[fill=cyan!30] (0,0) rectangle (1.5,1);
        \node at (0.75,0.5) {$45°$};
    \end{scope}
    
    % 90 градусов
    \begin{scope}[xshift=6cm, rotate=90]
        \draw[fill=cyan!30] (0,0) rectangle (1.5,1);
        \node at (0.75,0.5) {$90°$};
    \end{scope}
\end{tikzpicture}
\caption{Поворот изображения на разные углы}
\label{fig:angles}
\end{figure}

\subsubsection{Параметр scale}

Масштабирование изображения демонстрируется на рисунке \ref{fig:scaling}.

\begin{figure}[h]
\centering
\begin{tikzpicture}
    % scale=0.5
    \begin{scope}[scale=0.5]
        \draw[fill=purple!30] (0,0) rectangle (2,2);
        \node at (1,1) {\tiny scale=0.5};
    \end{scope}
    
    % scale=1.0
    \begin{scope}[xshift=2cm, scale=1.0]
        \draw[fill=purple!30] (0,0) rectangle (2,2);
        \node at (1,1) {\small scale=1.0};
    \end{scope}
    
    % scale=1.5
    \begin{scope}[xshift=5cm, scale=1.5]
        \draw[fill=purple!30] (0,0) rectangle (2,2);
        \node at (1,1) {scale=1.5};
    \end{scope}
\end{tikzpicture}
\caption{Различные значения масштаба}
\label{fig:scaling}
\end{figure}

\subsection{Разница между textwidth и linewidth}

\texttt{textwidth} — это полная ширина текстового блока на странице, в то время как \texttt{linewidth} — это текущая ширина строки, которая может изменяться внутри окружений (например, в minipage или в колонках).

\subsubsection{Демонстрация в обычном режиме}

В обычном однострочном режиме \texttt{textwidth} и \texttt{linewidth} имеют одинаковое значение.

\begin{figure}[h]
\centering
\begin{tikzpicture}
    \draw[fill=yellow!30] (0,0) rectangle (0.8\textwidth, 1);
    \node at (0.4\textwidth, 0.5) {width=0.8\textbackslash textwidth};
\end{tikzpicture}
\caption{Изображение с шириной относительно textwidth}
\label{fig:textwidth_demo}
\end{figure}

\begin{figure}[h]
\centering
\begin{tikzpicture}
    \draw[fill=green!30] (0,0) rectangle (0.8\linewidth, 1);
    \node at (0.4\linewidth, 0.5) {width=0.8\textbackslash linewidth};
\end{tikzpicture}
\caption{Изображение с шириной относительно linewidth}
\label{fig:linewidth_demo}
\end{figure}

\subsubsection{Демонстрация в minipage}

Внутри minipage \texttt{linewidth} изменяется, а \texttt{textwidth} остается прежним:

\begin{figure}[h]
\centering
\begin{minipage}{0.45\textwidth}
    \centering
    \begin{tikzpicture}[scale=0.6]
        \draw[fill=pink!40] (0,0) rectangle (4,2);
        \node at (2,1) {\small 0.8\textbackslash linewidth};
    \end{tikzpicture}
    \caption{В minipage linewidth меньше}
    \label{fig:minipage_line}
\end{minipage}
\hfill
\begin{minipage}{0.45\textwidth}
    \centering
    \small
    Здесь \texttt{linewidth} равен ширине minipage (0.45\textbackslash textwidth), в то время как \texttt{textwidth} по-прежнему равен полной ширине страницы.
\end{minipage}
\end{figure}

\subsection{Использование пакета lipsum для длинного текста}

Для демонстрации поведения плавающих объектов создадим документ с большим количеством текста.

\lipsum[1]

\begin{figure}[h]
\centering
\begin{tikzpicture}
    \draw[fill=blue!20] (0,0) rectangle (6,3);
    \node at (3,1.5) {\Large Спецификатор [h]};
\end{tikzpicture}
\caption{Плавающий объект со спецификатором h (here)}
\label{fig:float_h}
\end{figure}

\lipsum[2]

\begin{figure}[t]
\centering
\begin{tikzpicture}
    \draw[fill=red!20] (0,0) rectangle (6,3);
    \node at (3,1.5) {\Large Спецификатор [t]};
\end{tikzpicture}
\caption{Плавающий объект со спецификатором t (top)}
\label{fig:float_t}
\end{figure}

\lipsum[3]

\begin{figure}[b]
\centering
\begin{tikzpicture}
    \draw[fill=green!20] (0,0) rectangle (6,3);
    \node at (3,1.5) {\Large Спецификатор [b]};
\end{tikzpicture}
\caption{Плавающий объект со спецификатором b (bottom)}
\label{fig:float_b}
\end{figure}

\lipsum[4]

\subsection{Взаимодействие различных спецификаторов}

\subsubsection{Спецификатор [htbp]}

Это наиболее распространенная комбинация, дающая LaTeX максимальную свободу выбора:
\begin{figure}[htbp]
\centering
\begin{tikzpicture}
    \draw[fill=orange!20] (0,0) rectangle (6,2.5);
    \node at (3,1.25) {\Large Спецификатор [htbp]};
\end{tikzpicture}
\caption{Комбинированный спецификатор htbp}
\label{fig:float_htbp}
\end{figure}

\lipsum[5]

\subsubsection{Спецификатор [H] — строгое позиционирование}

Пакет \texttt{float} предоставляет спецификатор H для строгого размещения:

\begin{figure}[H]
\centering
\begin{tikzpicture}
    \draw[fill=purple!20] (0,0) rectangle (6,2.5);
    \node at (3,1.25) {\Large Спецификатор [H]};
\end{tikzpicture}
\caption{Строгое позиционирование с [H]}
\label{fig:float_H}
\end{figure}

Объект с спецификатором [H] всегда размещается строго в месте вставки в коде.

\lipsum[6]

\subsection{Работа с командами label и ref}

\subsubsection{Количество компиляций}

Для корректной работы перекрестных ссылок требуется \textbf{минимум две компиляции}:

\begin{enumerate}
    \item \textbf{Первая компиляция:} LaTeX записывает информацию о метках в файл .aux
    \item \textbf{Вторая компиляция:} LaTeX читает информацию из .aux и подставляет правильные номера
\end{enumerate}

Создадим несколько разделов для демонстрации:

\section{Тестовый раздел для ссылок}
\label{sec:test1}

Это раздел \ref{sec:test1}. Он содержит подраздел \ref{subsec:test1_1}.

\subsection{Тестовый подраздел}
\label{subsec:test1_1}

Ссылка на основной раздел: раздел \ref{sec:test1}.

\subsection{Еще один подраздел}
\label{subsec:test1_2}

Ссылки на предыдущие разделы: \ref{sec:test1}, \ref{subsec:test1_1}, \ref{subsec:test1_2}.

\subsubsection{Нумерованные списки}

\begin{enumerate}
    \item Первый элемент \label{item:1}
    \item Второй элемент \label{item:2}
    \item Третий элемент \label{item:3}
\end{enumerate}

Ссылка на элемент списка: пункт \ref{item:2}.

\subsection{Правильное размещение label относительно caption}

\subsubsection{Неправильное размещение (label перед caption)}

Если поместить \texttt{label} \textbf{перед} \texttt{caption}, ссылка будет указывать на номер раздела, а не на номер рисунка:

\begin{verbatim}
\begin{figure}[h]
\centering
[изображение]
\label{fig:wrong}  % НЕПРАВИЛЬНО!
\caption{Подпись}
\end{figure}
\end{verbatim}

\begin{figure}[h]
\centering
\begin{tikzpicture}
    \draw[fill=red!40] (0,0) rectangle (4,2);
    \node at (2,1) {\textbf{НЕПРАВИЛЬНО}};
\end{tikzpicture}
\label{fig:wrong_order}  % Это даст неправильную ссылку!
\caption{Пример неправильного размещения label}
\end{figure}

Ссылка на рисунок выше: \ref{fig:wrong_order} — может быть неправильной!

\subsubsection{Правильное размещение (label после caption)}

Правильный порядок: \texttt{caption} \textbf{затем} \texttt{label}:

\begin{verbatim}
\begin{figure}[h]
\centering
[изображение]
\caption{Подпись}
\label{fig:correct}  % ПРАВИЛЬНО!
\end{figure}
\end{verbatim}

\begin{figure}[h]
\centering
\begin{tikzpicture}
    \draw[fill=green!40] (0,0) rectangle (4,2);
    \node at (2,1) {\textbf{ПРАВИЛЬНО}};
\end{tikzpicture}
\caption{Пример правильного размещения label}
\label{fig:correct_order}
\end{figure}

Ссылка на рисунок выше: рисунок \ref{fig:correct_order} — корректная ссылка!

\subsection{Label для уравнений}

\subsubsection{Правильное размещение label в уравнении}

Для уравнений \texttt{label} должен быть \textbf{внутри} окружения equation:

\begin{equation}
E = mc^2
\label{eq:einstein}
\end{equation}

Ссылка на уравнение \ref{eq:einstein} работает корректно.

\subsubsection{Неправильное размещение (label после end{equation})}

Если поместить label \textbf{после} \texttt{end\{equation\}}, ссылка будет некорректной:

\begin{equation}
F = ma
\end{equation}
\label{eq:newton_wrong}  % НЕПРАВИЛЬНО!

Попытка сослаться на уравнение: \ref{eq:newton_wrong} — даст номер раздела, а не уравнения!

\subsubsection{Правильный пример}

\begin{equation}
\int_{a}^{b} f(x)dx = F(b) - F(a)
\label{eq:fundamental}
\end{equation}
Уравнение \ref{eq:fundamental} демонстрирует основную теорему математического анализа.

\section{Эксперименты с двухколоночным режимом}

Для демонстрации поведения изображений в режиме \texttt{twocolumn} создадим отдельный пример.

\subsection{Поведение textwidth и linewidth в twocolumn}

В режиме \texttt{twocolumn}:
\begin{itemize}
    \item \texttt{textwidth} — полная ширина страницы
    \item \texttt{linewidth} — ширина одной колонки
    \item \texttt{columnwidth} — ширина колонки (равна linewidth)
\end{itemize}

\subsection{Изображения на всю ширину страницы}

Для размещения изображения на всю ширину в режиме twocolumn используется окружение \texttt{figure*}:

\begin{verbatim}
\begin{figure*}[t]
\centering
\includegraphics[width=\textwidth]{image}
\caption{Изображение на всю ширину}
\end{figure*}
\end{verbatim}

\section{Дополнительные эксперименты}

\subsection{Обтекание текстом}

Пакет \texttt{wrapfig} позволяет создавать изображения, обтекаемые текстом:

\begin{wrapfigure}{r}{0.4\textwidth}
\centering
\begin{tikzpicture}[scale=0.8]
    \draw[fill=cyan!30] (0,0) circle (1.5);
    \node at (0,0) {\textbf{Wrap}};
\end{tikzpicture}
\caption{Обтекаемое изображение}
\label{fig:wrap}
\end{wrapfigure}

\lipsum[7]

\subsection{Множественные изображения}

Пакет \texttt{subcaption} позволяет размещать несколько изображений в одном float:

\begin{figure}[h]
\centering
\begin{subfigure}{0.3\textwidth}
    \centering
    \begin{tikzpicture}[scale=0.6]
        \draw[fill=red!30] (0,0) rectangle (2,2);
        \node at (1,1) {(a)};
    \end{tikzpicture}
    \caption{Первое}
    \label{fig:sub1}
\end{subfigure}
\hfill
\begin{subfigure}{0.3\textwidth}
    \centering
    \begin{tikzpicture}[scale=0.6]
        \draw[fill=green!30] (0,0) rectangle (2,2);
        \node at (1,1) {(b)};
    \end{tikzpicture}
    \caption{Второе}
    \label{fig:sub2}
\end{subfigure}
\hfill
\begin{subfigure}{0.3\textwidth}
    \centering
    \begin{tikzpicture}[scale=0.6]
        \draw[fill=blue!30] (0,0) rectangle (2,2);
        \node at (1,1) {(c)};
    \end{tikzpicture}
    \caption{Третье}
    \label{fig:sub3}
\end{subfigure}
\caption{Три изображения в одной figure}
\label{fig:multiple}
\end{figure}

Ссылки на подрисунки: \ref{fig:sub1}, \ref{fig:sub2}, \ref{fig:sub3}, и на всю группу: \ref{fig:multiple}.

\section{Результаты экспериментов}

\subsection{Параметры изображений}

В ходе экспериментов были получены следующие результаты:

\begin{table}[h]
\centering
\caption{Результаты экспериментов с параметрами}
\label{tab:results}
\begin{tabular}{|l|p{8cm}|}
\hline
\textbf{Параметр} & \textbf{Результат} \\
\hline
width & Позволяет задать ширину относительно textwidth или в абсолютных единицах \\
\hline
height & Задает высоту изображения, может нарушить пропорции \\
\hline
scale & Пропорционально масштабирует изображение \\
\hline
angle & Поворачивает изображение против часовой стрелки \\
\hline
keepaspectratio & Сохраняет пропорции при задании width и height \\
\hline
\end{tabular}
\end{table}

\subsection{Спецификаторы позиционирования}

Результаты тестирования спецификаторов:

\begin{enumerate}
    \item \textbf{[h]} — пытается разместить здесь, но может переместить
    \item \textbf{[t]} — размещает вверху страницы
    \item \textbf{[b]} — размещает внизу страницы
    \item \textbf{[p]} — создает отдельную страницу для floats
    \item \textbf{[H]} — строго в месте вставки (требует пакет float)
    \item \textbf{[htbp]} — дает LaTeX максимальную свободу выбора
\end{enumerate}

\subsection{Система перекрестных ссылок}

Основные выводы по работе с \texttt{label} и \texttt{ref}:

\begin{itemize}
    \item Требуется \textbf{минимум 2 компиляции} для корректной работы ссылок
    \item Для сложных документов может потребоваться 3-4 компиляции
    \item \texttt{label} должен идти \textbf{после} \texttt{caption} для рисунков и таблиц
    \item \texttt{label} должен быть \textbf{внутри} окружения equation
    \item Неправильное размещение приводит к ссылкам на номер раздела вместо объекта
\end{itemize}

\section{Выводы}

В ходе выполнения лабораторной работы были достигнуты следующие результаты:

\begin{enumerate}
    \item Освоены навыки создания векторной графики с помощью пакета TikZ
    \item Изучены все основные параметры команды \texttt{includegraphics}
    \item Выявлена разница между \texttt{textwidth} и \texttt{linewidth}
    \item Протестировано поведение изображений в режиме twocolumn
    \item Исследованы все спецификаторы позиционирования плавающих объектов
    \item Определено необходимое количество компиляций для работы ссылок
    \item Установлено правильное размещение \texttt{label} относительно \texttt{caption}
    \item Выявлены последствия неправильного размещения \texttt{label}
\end{enumerate}

\subsection{Практические рекомендации}

На основе проведенных экспериментов можно дать следующие рекомендации:

\begin{itemize}
    \item Используйте относительные размеры (0.8\textbackslash textwidth) вместо абсолютных
    \item Для сохранения пропорций используйте параметр \texttt{keepaspectratio=true}
    \item Предпочитайте спецификатор [htbp] для большинства случаев
    \item Используйте [H] только когда строгое позиционирование критично
    \item Всегда размещайте \texttt{label} после \texttt{caption}
    \item Компилируйте документ минимум дважды после добавления ссылок
    \item В режиме twocolumn используйте \texttt{figure*} для широких изображений
\end{itemize}

\section{Список использованных источников}

\begin{enumerate}
    \item The LaTeX Graphics Companion — \url{https://www.latex-project.org/}
    \item TikZ \& PGF Manual — \url{https://tikz.dev/}
    \item Документация пакета graphicx
    \item Документация пакета float
    \item Wikibooks: LaTeX/Floats, Figures and Captions
\end{enumerate}

\newpage

\section*{Приложение А. Пример кода для создания изображений}
\addcontentsline{toc}{section}{Приложение А. Пример кода}

\subsection*{Код для создания геометрической композиции}

\begin{verbatim}
\begin{tikzpicture}
    \draw[fill=blue!30, draw=blue!50!black, line width=2pt] 
        (0,0) rectangle (4,3);
    \draw[fill=red!30, draw=red!50!black, line width=2pt] 
        (2,1.5) circle (1cm);
    \draw[fill=green!30, draw=green!50!black, line width=2pt] 
        (1,2.5) -- (3,2.5) -- (2,4) -- cycle;
    \node at (2,0.5) {\Large\textbf{TikZ}};
\end{tikzpicture}
\end{verbatim}

\subsection*{Код для графика функции}

\begin{verbatim}
\begin{tikzpicture}[scale=1.2]
    \draw[->, thick] (-0.5,0) -- (5,0) node[right] {$x$};
    \draw[->, thick] (0,-0.5) -- (0,4) node[above] {$y$};
    \draw[domain=0:4.5, smooth, variable=\x, blue, line width=1.5pt] 
        plot ({\x}, {0.5*\x*\x - 0.3*\x + 0.5});
    \fill[red] (1,0.7) circle (2pt) node[above right] {$A$};
    \draw[gray!30, very thin] (0,0) grid (4.5,4);
\end{tikzpicture}
\end{verbatim}

\end{document}
